% ======================================================================
%  CAPÍTULO 2 — FUNDAMENTAÇÃO TEÓRICA
%
%  Os valores numéricos do ArduPilot e do SITL citados aqui foram
%  conferidos no código-fonte (commit ce9c2d01a4) e nos params.parm
%  gravados pelos ensaios em ensaios/resultados/.
% ======================================================================
\chapter{Fundamentação teórica}
\label{cap:fundamentacao}

Este capítulo reúne a base teórica sobre a qual o trabalho se apoia. A
Seção~\ref{sec:atitude} define os referenciais e as representações de atitude;
as Seções~\ref{sec:modelo} e~\ref{sec:linearizacao} derivam o modelo dinâmico
do quadrirrotor e sua aproximação linear em torno do voo pairado. A
Seção~\ref{sec:ardupilot} descreve como o ArduPilot estima a atitude e a
controla, e onde os controladores propostos se inserem nessa cadeia. As
Seções~\ref{sec:tecnicas} a~\ref{sec:identificacao} tratam das técnicas de
controle, dos modelos de perturbação e da identificação no domínio da
frequência, e a Seção~\ref{sec:sitl} descreve o simulador em que os ensaios
são executados.

% ----------------------------------------------------------------------
\section{Referenciais e representação da atitude}
\label{sec:atitude}

Adotam-se dois referenciais, ambos dextrógiros e consistentes com a convenção
do ArduPilot. O referencial inercial local $\{\mathcal{N}\}$ é o NED
(\textit{North-East-Down}), com origem fixa na superfície e eixos apontando
para o norte, para o leste e para baixo; a hipótese de Terra plana e sem
rotação é adequada às distâncias e às durações envolvidas~\cite[Cap.~3]{beard2012}. O
referencial do corpo $\{\mathcal{B}\}$ é o FRD (\textit{Forward-Right-Down}),
com origem no centro de massa e eixos $x_b$ para a frente, $y_b$ para a
direita e $z_b$ para baixo (Figura~\ref{fig:quadx}).

\subsection{Ângulos de Euler}

A orientação de $\{\mathcal{B}\}$ em relação a $\{\mathcal{N}\}$ pode ser
descrita pelos ângulos de Euler na sequência $Z$-$Y$-$X$: guinada $\psi$ em
torno de $z$, arfagem $\theta$ em torno do $y$ intermediário e rolagem $\phi$
em torno de $x_b$. A matriz de rotação que leva vetores do corpo para o
referencial inercial é~\cite{beard2008,sabatino2015,beard2012}
\begin{equation}
\matriz{R} = \matriz{R}_z(\psi)\,\matriz{R}_y(\theta)\,\matriz{R}_x(\phi) =
\begin{bmatrix}
c_\theta c_\psi & s_\phi s_\theta c_\psi - c_\phi s_\psi & c_\phi s_\theta c_\psi + s_\phi s_\psi \\
c_\theta s_\psi & s_\phi s_\theta s_\psi + c_\phi c_\psi & c_\phi s_\theta s_\psi - s_\phi c_\psi \\
-s_\theta       & s_\phi c_\theta                         & c_\phi c_\theta
\end{bmatrix},
\label{eq:rotacao}
\end{equation}
em que $c_\alpha = \cos\alpha$ e $s_\alpha = \sin\alpha$. A velocidade
angular do corpo, $\vetor{\omega} = [p\ q\ r]^{\mathsf{T}}$, é a grandeza
medida pelos giroscópios e não coincide com as derivadas dos ângulos de Euler.
A relação entre ambas é~\cite{beard2012}
\begin{equation}
\begin{bmatrix} \dot\phi \\ \dot\theta \\ \dot\psi \end{bmatrix} =
\begin{bmatrix}
1 & s_\phi t_\theta & c_\phi t_\theta \\
0 & c_\phi          & -s_\phi \\
0 & s_\phi/c_\theta & c_\phi/c_\theta
\end{bmatrix}
\vetor{\omega},
\qquad t_\theta = \tan\theta,
\label{eq:cinematica-euler}
\end{equation}
singular em $\theta = \pm\pi/2$. Os ângulos de Euler são a representação mais
intuitiva para especificar referências e apresentar resultados, e é nela que
as métricas deste trabalho são calculadas. Não são, contudo, a representação
usada internamente pelo firmware.

\subsection{Quatérnios, matriz de cossenos diretores e vetor de rotação}

O ArduPilot representa a atitude por quatérnios unitários
$\vetor{q} = [q_0\ q_1\ q_2\ q_3]^{\mathsf{T}} =
[\cos(\alpha/2)\ \ \vetor{e}^{\mathsf{T}}\sin(\alpha/2)]^{\mathsf{T}}$, que
descrevem uma rotação de ângulo $\alpha$ em torno do eixo unitário
$\vetor{e}$, sem singularidades. A biblioteca matemática do firmware segue a
convenção de Hamilton para o produto de quatérnios, e seus testes unitários
conferem a conversão de quatérnio para matriz de rotação contra o exemplo
publicado por \citeonline{sommer2018}\footnote{O código do ArduPilot cita a
versão do artigo publicada no arXiv,
\url{https://doi.org/10.48550/arXiv.1801.07478}.}, que discutem as
inconsistências causadas pela convenção alternativa. Nessa convenção, a
cinemática de atitude é~\cite{sola2017,hager2026}
\begin{equation}
\dot{\vetor{q}} = \frac{1}{2}\,\vetor{q} \otimes
\begin{bmatrix} 0 \\ \vetor{\omega} \end{bmatrix},
\label{eq:cinematica-quaternio}
\end{equation}
com $\otimes$ o produto de Hamilton e $\vetor{q}$ a rotação do corpo para o
referencial inercial.

Duas outras representações aparecem no código do ArduPilot e na literatura
que ele referencia. A primeira é a própria matriz $\matriz{R}$, ou matriz de
cossenos diretores (DCM), integrada diretamente por
$\dot{\matriz{R}} = \matriz{R}\,[\vetor{\omega}]_\times$ e periodicamente
re-ortogonalizada: é a base do estimador \codigo{AP\_AHRS\_DCM}, derivado do
trabalho de \citeonline{premerlani2009}. A segunda é o vetor de rotação
$\delta\vetor{\theta} = \alpha\,\vetor{e}$, uma parametrização mínima, de três
componentes, adequada a representar \emph{erros} pequenos de atitude; é a
formulação de erro adotada nas derivações simbólicas de filtros de Kalman que
acompanham o código, conforme \citeonline{pittelkau2003}.

A distinção entre a representação interna e a de apresentação tem
consequência prática. O firmware modela a referência de atitude como uma
rotação suave em quatérnio; quando os três eixos são comandados ao mesmo
tempo, os ângulos de Euler dessa referência deixam de evoluir de forma
independente, e a referência expressa em rolagem e arfagem pode ultrapassar
transitoriamente o valor final comandado (Seção~\ref{sec:res-degrau}), sem
que o controlador tenha participação nisso. Por esse motivo, nos ensaios com os três eixos
simultâneos, a comparação entre controladores se faz pelo erro de
rastreamento em relação à referência efetivamente modelada, e não apenas pelo
sobressinal em relação ao degrau nominal.

% ----------------------------------------------------------------------
\section{Modelo dinâmico do quadrirrotor}
\label{sec:modelo}

\subsection{Princípio de funcionamento e configuração X}

O quadrirrotor possui quatro rotores de passo fixo, montados nas extremidades
de braços de comprimento $\ell$ e girando aos pares em sentidos opostos, de
modo que os torques de reação se cancelem em equilíbrio. Na configuração X,
adotada neste trabalho, os braços formam ângulos de \ang{45} com o eixo
longitudinal. A Figura~\ref{fig:quadx} mostra a numeração e o sentido de
rotação dos motores segundo a convenção do ArduPilot.

\begin{figure}[htbp]
\centering
\caption{Quadrirrotor em configuração X, visto de cima, com a numeração e o
sentido de rotação dos motores adotados pelo ArduPilot.}
\label{fig:quadx}
\begin{tikzpicture}[
    font=\small,
    seta/.style={-{Stealth[length=2.2mm]}},
    giro/.style={-{Stealth[length=2mm]}, thick},
    cota/.style={{Stealth[length=1.6mm]}-{Stealth[length=1.6mm]}},
]
    % braços
    \draw[line width=2.2pt, gray!70] (-2,-2) -- (2,2);
    \draw[line width=2.2pt, gray!70] (-2,2) -- (2,-2);
    \fill[gray!40] (0,0) circle (0.42);
    % rotores: posição na página (x = direita = y_b, y = cima = x_b)
    \foreach \x/\y/\n/\s in {2/2/1/CCW, -2/-2/2/CCW, -2/2/3/CW, 2/-2/4/CW} {
        \fill[white] (\x,\y) circle (0.72);
        \draw[thick] (\x,\y) circle (0.72);
        \node at (\x,\y+0.1) {\textbf{M\n}};
        \node[font=\scriptsize] at (\x,\y-0.3) {\s};
    }
    % sentido de giro visto de cima (anti-horário = CCW)
    \foreach \x/\y in {2/2, -2/-2}
        \draw[giro, blue!70!black] ($(\x,\y)+(-30:0.95)$)
            arc[start angle=-30, end angle=210, radius=0.95];
    \foreach \x/\y in {-2/2, 2/-2}
        \draw[giro, red!70!black] ($(\x,\y)+(210:0.95)$)
            arc[start angle=210, end angle=-30, radius=0.95];
    % eixos do corpo; z_b entra na página
    \draw[seta, very thick] (0,0.42) -- (0,1.45) node[above] {$x_b$ (frente)};
    \draw[seta, very thick] (0.42,0) -- (1.45,0) node[right] {$y_b$};
    \draw[thick, fill=white] (0,0) circle (0.13);
    \draw[thick] (-0.09,-0.09) -- (0.09,0.09) (-0.09,0.09) -- (0.09,-0.09);
    \node[left] at (-0.45,0) {$z_b$};
    % braço e alavanca
    \node[above left=-2pt, font=\footnotesize] at (1.0,1.0) {$\ell$};
    \draw[densely dotted] (0,-0.45) -- (0,-3.15);
    \draw[densely dotted] (2,-2.72) -- (2,-3.15);
    \draw[cota] (0,-3.0) -- (2,-3.0) node[midway, below] {$\ell_{xy}$};
\end{tikzpicture}
\fonte{Elaborada pelo autor.}
\end{figure}

A rolagem é produzida pela diferença de empuxo entre os lados direito e
esquerdo, e a arfagem, pela diferença entre a frente e a traseira. A guinada
não tem atuador próprio: resulta do desequilíbrio entre os torques de reação
dos pares que giram em sentidos opostos. Como esse torque de reação é muito
menor que o momento que o mesmo empuxo produz com o braço, a autoridade de
controle em guinada é muito inferior à de rolagem e arfagem --- característica
quantificada na Seção~\ref{sec:sitl} para a aeronave simulada.

\subsection{Hipóteses de modelagem}

O modelo é derivado sob as hipóteses usuais da
literatura~\cite{bouabdallah2007,beard2008,sabatino2015}:
\begin{alineas}
    \item a aeronave é um corpo rígido de massa constante, simétrico em relação
          aos planos $x_b z_b$ e $y_b z_b$, de modo que o tensor de inércia é
          diagonal, $\matriz{J} = \operatorname{diag}(J_{xx}, J_{yy}, J_{zz})$;
    \item o centro de massa coincide com a origem de $\{\mathcal{B}\}$ e está
          no plano dos rotores;
    \item o empuxo e o torque de reação de cada rotor são proporcionais ao
          quadrado da sua velocidade de rotação $\Omega_i$;
    \item a dinâmica dos motores e dos controladores eletrônicos de
          velocidade, o arrasto do corpo, o \textit{blade flapping} e o efeito
          solo são desprezados no modelo nominal e tratados como fontes
          conhecidas de discrepância.
\end{alineas}

\subsection{Equações de Newton-Euler}

A dinâmica de translação, escrita no referencial inercial, é
\begin{equation}
\dot{\vetor{p}} = \vetor{v}, \qquad
m\,\dot{\vetor{v}} = m g\,\vetor{e}_3 - T\,\matriz{R}\,\vetor{e}_3 + \vetor{f}_d,
\label{eq:translacao}
\end{equation}
em que $\vetor{p}$ e $\vetor{v}$ são a posição e a velocidade do centro de
massa, $\vetor{e}_3 = [0\ 0\ 1]^{\mathsf{T}}$, $T$ é o empuxo total, que atua
ao longo de $-z_b$, e $\vetor{f}_d$ reúne as forças de perturbação, entre elas
a força aerodinâmica devida ao vento. A dinâmica de rotação, escrita no
referencial do corpo, é dada pela equação de Euler~\cite{beard2008,beard2012}:
\begin{equation}
\matriz{J}\,\dot{\vetor{\omega}} =
-\vetor{\omega}\times\matriz{J}\,\vetor{\omega} + \vetor{\tau} + \vetor{\tau}_d,
\label{eq:rotacao-euler}
\end{equation}
em que $\vetor{\tau} = [\tau_x\ \tau_y\ \tau_z]^{\mathsf{T}}$ é o torque
produzido pelos rotores e $\vetor{\tau}_d$ é um torque de perturbação aditivo.
É por $\vetor{\tau}_d$ que entram, neste trabalho, os efeitos de momento não
modelados --- assimetrias, rajadas que atingem os rotores de forma desigual,
perda de eficiência de um rotor --- e as perturbações de torque usadas nos
ensaios (Seção~\ref{sec:perturbacoes}).

O conjunto formado por \eqref{eq:translacao}, \eqref{eq:rotacao-euler} e pela
cinemática \eqref{eq:cinematica-euler} ou \eqref{eq:cinematica-quaternio}
define um sistema de doze estados,
$\vetor{x} = [\vetor{p}^{\mathsf{T}}\ \vetor{v}^{\mathsf{T}}\
\phi\ \theta\ \psi\ \vetor{\omega}^{\mathsf{T}}]^{\mathsf{T}}$, com quatro
entradas, $\vetor{u} = [T\ \tau_x\ \tau_y\ \tau_z]^{\mathsf{T}}$. A
subatuação fica evidente: as acelerações horizontais não podem ser comandadas
diretamente, apenas por meio da inclinação do vetor de empuxo, o que impõe a
estrutura em cascata em que a atitude é a malha interna.

\subsection{Modelo dos rotores e matriz de alocação}

Cada rotor $i$ produz um empuxo $T_i = k_T\,\Omega_i^2$ ao longo de $-z_b$ e
um torque de reação $Q_i = k_Q\,\Omega_i^2 = \kappa\,T_i$ em torno de $z_b$,
com $\kappa = k_Q/k_T$. O torque de reação tem sentido oposto ao da rotação do
rotor: um rotor que gira no sentido anti-horário visto de cima (CCW) produz
sobre o corpo um torque no sentido horário, isto é, positivo em torno de
$z_b$. Com o motor $i$ na posição $(x_i, y_i)$ do plano dos rotores, o torque
do empuxo é $(x_i, y_i, 0)\times(0, 0, -T_i) = (-y_i T_i,\ x_i T_i,\ 0)$. Na
configuração X, $|x_i| = |y_i| = \ell_{xy} = \ell/\sqrt{2}$, e a numeração da
Figura~\ref{fig:quadx} leva à matriz de alocação
\begin{equation}
\begin{bmatrix} T \\ \tau_x \\ \tau_y \\ \tau_z \end{bmatrix} =
\underbrace{\begin{bmatrix}
1          & 1          & 1          & 1 \\
-\ell_{xy} & \ell_{xy}  & \ell_{xy}  & -\ell_{xy} \\
\ell_{xy}  & -\ell_{xy} & \ell_{xy}  & -\ell_{xy} \\
\kappa     & \kappa     & -\kappa    & -\kappa
\end{bmatrix}}_{\matriz{M}}
\begin{bmatrix} T_1 \\ T_2 \\ T_3 \\ T_4 \end{bmatrix}.
\label{eq:alocacao}
\end{equation}
A matriz $\matriz{M}$ é invertível, e $\matriz{M}^{-1}$ fornece os empuxos
individuais que realizam um par empuxo--torque desejado, desde que nenhum
deles saia da faixa disponível. A comparação entre a segunda e a quarta linhas
explica a fraca autoridade em guinada: um mesmo padrão de diferença de empuxo
produz torques na razão $\ell_{xy}/\kappa$ entre rolagem e guinada, da ordem
de 25 na aeronave simulada; como o torque de reação do SITL cresce com o
comando, a razão que vale para variações de empuxo é $\ell_{xy}/\kappa'$, com
$\kappa' = \mathrm{d}Q_i/\mathrm{d}T_i$, cerca de 15
(Seção~\ref{sec:forcas-momentos}).

% ----------------------------------------------------------------------
\section{Linearização em torno do voo pairado}
\label{sec:linearizacao}

O ponto de equilíbrio de interesse é o voo pairado: $\vetor{v} = \vetor{0}$,
$\phi = \theta = 0$, $\psi = \psi_0$, $\vetor{\omega} = \vetor{0}$, com
$T = mg$ e $\vetor{\tau} = \vetor{0}$. Tomando, sem perda de generalidade,
$\psi_0 = 0$ e desprezando produtos de pequenas variações, obtém-se de
\eqref{eq:rotacao}--\eqref{eq:rotacao-euler}
\begin{equation}
\delta\ddot{x} = -g\,\delta\theta, \quad
\delta\ddot{y} = g\,\delta\phi, \quad
\delta\ddot{z} = -\frac{\delta T}{m}, \quad
\ddot\phi = \frac{\tau_x}{J_{xx}}, \quad
\ddot\theta = \frac{\tau_y}{J_{yy}}, \quad
\ddot\psi = \frac{\tau_z}{J_{zz}}.
\label{eq:linear}
\end{equation}
O termo giroscópico $\vetor{\omega}\times\matriz{J}\vetor{\omega}$ e o
acoplamento da cinemática \eqref{eq:cinematica-euler} são de segunda ordem e
desaparecem; o sistema se decompõe em subsistemas de altitude, de atitude e de
posição horizontal, este último alimentado pela atitude. Essa estrutura
fundamenta tanto a arquitetura em cascata do firmware quanto o projeto do
controlador por eixo.

Cada eixo de atitude é, em primeira aproximação, um duplo
integrador~\cite{rao2001}. Para a
rolagem, com estado $\vetor{x}_\phi = [\phi\ \ p]^{\mathsf{T}}$,
\begin{equation}
\dot{\vetor{x}}_\phi =
\underbrace{\begin{bmatrix} 0 & 1 \\ 0 & 0 \end{bmatrix}}_{\matriz{A}}
\vetor{x}_\phi +
\underbrace{\begin{bmatrix} 0 \\ 1/J_{xx} \end{bmatrix}}_{\matriz{B}}
(\tau_x + \tau_{d,x}),
\qquad
\frac{\Phi(s)}{\tau_x(s)} = \frac{1}{J_{xx}\,s^2},
\label{eq:duplo-integrador}
\end{equation}
e analogamente para arfagem e guinada. O par $(\matriz{A}, \matriz{B})$ é
controlável, pois $[\matriz{B}\ \ \matriz{A}\matriz{B}]$ tem posto 2, e o par
$(\matriz{A}, \matriz{C})$ é observável já com
$\matriz{C} = [1\ \ 0]$~\cite{ogata2010}; na
prática, os dois estados são disponíveis, o ângulo pelo estimador e a
velocidade angular pelo giroscópio. Os dois autovalores na origem, num bloco
de Jordan, tornam o sistema instável em malha aberta. Em malha fechada, os
integradores da planta não ajudam a rejeitar um torque de perturbação
constante, porque ele entra antes deles: com realimentação proporcional dos
estados, $\tau_x = -J_{xx}\,(k_\theta\,\phi + k_\omega\,p)$, o ângulo se
acomoda num erro $\tau_{d,x}/(J_{xx}k_\theta)$. Erro nulo em regime exige um
integrador no controlador ou uma estimativa da perturbação, como prevê o
princípio do modelo interno~\cite{francis1976}.

Como o controlador é executado a cada $T_s$ segundos e o torque comandado é
mantido constante entre execuções, o modelo de projeto é a discretização por
segurador de ordem zero (ZOH)~\cite{franklin1998}:
\begin{equation}
\vetor{x}_{\phi,k+1} =
\begin{bmatrix} 1 & T_s \\ 0 & 1 \end{bmatrix} \vetor{x}_{\phi,k} +
\begin{bmatrix} T_s^2/(2J_{xx}) \\ T_s/J_{xx} \end{bmatrix} \tau_{x,k},
\label{eq:discreto}
\end{equation}
com $T_s = \SI{2.5}{\milli\second}$ para a taxa de \SI{400}{\hertz} usada
neste trabalho. O duplo integrador é o modelo \emph{nominal}; a planta efetiva
vista pelo controlador inclui ainda o filtro passa-baixas dos giroscópios, os
atrasos de amostragem e a saturação da alocação e, numa aeronave real, a
dinâmica dos motores, ausente no SITL (Seção~\ref{sec:motor-sitl}), que
acrescentam atraso de fase e limitam a banda passante alcançável. É essa
diferença que a identificação da Seção~\ref{sec:identificacao} quantifica.

% ----------------------------------------------------------------------
\section{Estimação e controle de atitude no ArduPilot}
\label{sec:ardupilot}

\subsection{Laço de execução}

No ArduCopter, a malha de atitude é executada no laço rápido do escalonador,
na taxa definida pelo parâmetro \codigo{SCHED\_LOOP\_RATE}, de
\SI{400}{\hertz} neste trabalho. A cada ciclo, na ordem definida em
\codigo{ArduCopter/Copter.cpp}~\cite{ardupilot2026}, o firmware lê os sensores
inerciais, executa o controlador de velocidade angular, executa o controlador
customizado, quando habilitado, envia os comandos aos motores e só então
atualiza o estimador de atitude e o modo de voo, cujo controlador de ângulo
produz o alvo de velocidade angular consumido no ciclo seguinte. A ordem
minimiza a latência entre a amostra do giroscópio e o comando dos motores. A
Figura~\ref{fig:cascata} resume a cadeia.

\begin{figure}[htbp]
\centering
\caption{Cadeia de controle de atitude do ArduCopter e ponto de inserção do
\codigo{AC\_CustomControl}. O modelador de referência fornece também a
velocidade angular de alimentação direta $\vetor{\omega}_{ff}$, e o
controlador customizado recebe as mesmas referências e estimativas que o
nativo.}
\label{fig:cascata}
\begin{tikzpicture}[
    font=\scriptsize,
    bloco/.style={draw, rectangle, align=center, minimum height=1.15cm,
                  text width=1.75cm, inner sep=2pt},
    seta/.style={-{Stealth[length=1.8mm]}},
    node distance=0.8cm,
]
    \node[bloco] (ref) {Modo de voo\\(referência)};
    \node[bloco, right=of ref] (mod) {Modelador de referência};
    \node[bloco, right=of mod] (ang) {Controlador de ângulo};
    \node[bloco, right=of ang] (taxa) {PID de velocidade angular};
    \node[bloco, right=of taxa] (mix) {Alocação de empuxo};
    \node[bloco, right=of mix] (planta) {Motores e aeronave};
    \node[bloco, below=0.9cm of taxa] (est) {Estimação\\(INS e EKF3)};
    \node[bloco, dashed, text width=3.3cm, minimum height=0.8cm]
        (cc) at ($(ang)!0.5!(taxa) + (0,1.55)$)
        {\codigo{AC\_CustomControl}\\(opcional, por eixo)};

    \draw[seta] (ref) -- (mod);
    \draw[seta] (mod) -- node[above] {$\vetor{q}^*$} (ang);
    \draw[seta] (ang) -- node[above] {$\vetor{\omega}^*$} (taxa);
    \draw[seta] (taxa) -- node[above] {$\vetor{u}$} (mix);
    \draw[seta] (mix) -- node[above] {PWM} (planta);
    \draw[seta] (planta.south) |- (est.east);
    \draw[seta] (est.west) -| node[pos=0.8, left] {$\hat{\vetor{q}}$} (ang.south);
    \draw[seta] (est.north) -- node[right] {$\vetor{\omega}$} (taxa.south);
    \draw[seta, dashed] (mod.north) |- (cc.west);
    \draw[seta, dashed] (cc.east) -| node[pos=0.75, right] {$\vetor{u}_{cc}$} (mix.north);
\end{tikzpicture}
\fonte{Elaborada pelo autor com base em \citeonline{ardupilot2026}.}
\end{figure}

\subsection{Estimação de atitude}

Os sensores inerciais são integrados em incrementos de ângulo e de
velocidade. Como rotações finitas não comutam, a simples soma dos incrementos
do giroscópio acumula erro quando a aeronave oscila em mais de um eixo ao
mesmo tempo (efeito de \textit{coning}); o firmware aplica a correção de
\textit{coning} descrita por Tian et al.\ (2010, p.~26) \textit{apud}
\citeonline{ardupilot2026} antes de entregar os
incrementos ao estimador.

O estimador padrão é o EKF3, um filtro de Kalman estendido que funde os
sensores inerciais com GPS, barômetro e magnetômetro e estende o vetor de
estados com as polarizações dos sensores inerciais, o campo magnético e o
vento~\cite{ardupilotdev,ardupilot2026}. O EKF3 calcula ainda, por um filtro
complementar de terceira ordem baseado em \citeonline{widnall1980}, que
combina a aceleração vertical com a altura estimada, uma derivada da altura
cinematicamente consistente com ela; as malhas de altitude só a usam no lugar
da velocidade vertical do EKF quando esta não está disponível ou sob vibração
elevada. Como alternativa de contingência, o firmware mantém o
estimador \codigo{AP\_AHRS\_DCM}, um filtro complementar sobre a matriz de
cossenos diretores~\cite{premerlani2009}. Em todos os casos, a malha de
atitude recebe a atitude estimada $\hat{\vetor{q}}$ e a velocidade angular
$\vetor{\omega}$ medida pelo giroscópio após um filtro passa-baixas
(\SI{20}{\hertz} neste trabalho, parâmetro \codigo{INS\_GYRO\_FILTER}).

\subsection{Modelador de referência}

A referência que chega do modo de voo --- do piloto, no modo estabilizado, ou
das malhas externas, nos modos automáticos --- não é aplicada diretamente ao
controlador de ângulo. Ela passa por um modelador que gera uma atitude alvo
$\vetor{q}^*$ cinematicamente consistente, com aceleração angular limitada
pelos parâmetros \codigo{ATC\_ACC\_R\_MAX}, \codigo{ATC\_ACC\_P\_MAX} e
\codigo{ATC\_ACC\_Y\_MAX} e suavizada por uma constante de tempo
(\codigo{ATC\_INPUT\_TC} em rolagem e arfagem e \codigo{PILOT\_Y\_RATE\_TC}
em guinada). O modelador, habilitado por \codigo{ATC\_RATE\_FF\_ENAB} (padrão
1; com 0, a atitude alvo recebe a referência sem limitação de aceleração),
produz também a velocidade angular correspondente à trajetória do alvo,
$\vetor{\omega}_{ff}$, que é somada à saída do controlador de ângulo como
alimentação direta~\cite{ardupilot2026}. Um degrau
comandado ao firmware, portanto, não chega ao controlador como degrau, e sim
como uma transição de aceleração limitada; os ensaios deste trabalho medem a
resposta a essa referência modelada, que é a mesma para todos os
controladores.

\subsection{Controlador de ângulo}

O erro de atitude é calculado em quatérnio e decomposto em duas rotações
sucessivas: a que alinha o vetor de empuxo estimado ao alvo, que define os
erros de rolagem e arfagem, e a rotação residual em torno do eixo do empuxo,
que define o erro de guinada. A decomposição permite priorizar o vetor de
empuxo: quando o erro de inclinação supera \ang{30}, a correção de guinada é
atenuada, e o erro de guinada é limitado ao valor que saturaria a saída em
regime~\cite{ardupilot2026}.

Cada componente $e$ do erro é convertida em velocidade angular alvo por um
controlador proporcional com limite de aceleração, chamado no código de
controlador de raiz quadrada:
\begin{equation}
\omega^*(e) =
\begin{cases}
k\,e, & |e| \le e_{\mathrm{lin}}, \\[2pt]
\operatorname{sgn}(e)\sqrt{2\,a_{\max}\left(|e| - e_{\mathrm{lin}}/2\right)},
   & |e| > e_{\mathrm{lin}},
\end{cases}
\qquad e_{\mathrm{lin}} = \frac{a_{\max}}{k^2},
\label{eq:sqrt-controller}
\end{equation}
em que $k$ é o ganho \codigo{ATC\_ANG\_*\_P} e $a_{\max}$ é metade da
aceleração máxima do eixo, \codigo{ATC\_ACC\_*\_MAX}/2, limitada à faixa de
40 a \SI{720}{\degree\per\second\squared} em rolagem e arfagem e à de 10 a
\SI{120}{\degree\per\second\squared} em guinada. Para erros pequenos, a lei é linear; para erros
grandes, a velocidade angular cresce com a raiz do erro, o que garante que a
aeronave consiga desacelerar até o alvo sem exceder $a_{\max}$. As duas
expressões e suas derivadas coincidem em $|e| = e_{\mathrm{lin}}$. O resultado
é limitado às velocidades angulares máximas configuradas e somado à
alimentação direta $\vetor{\omega}_{ff}$, expressa no corpo, formando o alvo
$\vetor{\omega}^*$ do controlador de velocidade angular.

\subsection{Controlador de velocidade angular}

Cada eixo tem um controlador PID (\codigo{AC\_PID}) sobre o erro de
velocidade angular, cuja saída é um esforço normalizado em $[-1, 1]$, fração
do torque máximo disponível no eixo:
\begin{equation}
u = k_P\,e_f + \int_0^t k_I\,e_f\,\mathrm{d}t' + k_D\,\frac{\mathrm{d}e_f}{\mathrm{d}t}
    + k_{FF}\,\omega^*_f,
\qquad e_f = F_E\{\omega^*_f - \omega\},\quad \omega^*_f = F_T\{\omega^*\},
\label{eq:pid-taxa}
\end{equation}
em que $F_T$ e $F_E$ são filtros passa-baixas de primeira ordem sobre o alvo e
sobre o erro (\codigo{FLTT} e \codigo{FLTE}), e o termo derivativo tem o seu
próprio filtro (\codigo{FLTD}). O integrador é limitado a \codigo{IMAX} e,
quando a biblioteca de motores sinaliza saturação no eixo, só pode diminuir
em módulo, como proteção contra \textit{windup}~\cite[Cap.~10]{astrom2008}. A Tabela~\ref{tab:ganhos}
reúne os ganhos usados neste trabalho.

\begin{table}[htbp]
\centering
\caption{Ganhos do controlador de atitude nativo na linha de base
(\codigo{nativo.yaml}).}
\label{tab:ganhos}
\begin{tabular}{lccc}
\toprule
Parâmetro & Rolagem & Arfagem & Guinada \\
\midrule
Ganho de ângulo $k$ (\si{\per\second})      & 26,7     & 25,4     & 2,29 \\
$k_P$ de velocidade angular                 & 0,0862   & 0,0862   & 0,608 \\
$k_I$ de velocidade angular                 & 0,0862   & 0,0862   & 0,0608 \\
$k_D$ de velocidade angular                 & 0,000375 & 0,000375 & 0 \\
Limite do integrador (\codigo{IMAX})        & 0,5      & 0,5      & 0,5 \\
Filtro do alvo, \codigo{FLTT} (\si{\hertz}) & 20       & 20       & 20 \\
Filtro do erro, \codigo{FLTE} (\si{\hertz}) & ---      & ---      & 1 \\
Filtro do derivativo, \codigo{FLTD} (\si{\hertz}) & 20 & 20       & 20 \\
Aceleração máxima (\si{\degree\per\second\squared}) & 5857 & 5703 & 115 \\
\bottomrule
\end{tabular}


\fonte{\citeonline{ardupilot2026}; parâmetros gravados nos ensaios.}
\nota{\raggedright Ganhos, filtro do erro de guinada e acelerações máximas
obtidos pelo AutoTune do ArduPilot na própria aeronave simulada
(Seção~\ref{sec:autotune}); filtros do alvo e do derivativo e limite do
integrador, que o AutoTune não ajusta, nos padrões compilados. ``---'' indica
filtro desabilitado. As variantes \codigo{nativo\_lqi} e
\codigo{nativo\_filtro10} voam o mesmo controlador com outros ganhos
(Seção~\ref{sec:controladores}).}
\end{table}

\subsection{Alocação de empuxo}

Os esforços normalizados de rolagem, arfagem e guinada e o comando de empuxo
coletivo são combinados pela biblioteca \codigo{AP\_MotorsMatrix}, que
implementa uma versão normalizada de $\matriz{M}^{-1}$ em
\eqref{eq:alocacao}: cada motor recebe o empuxo coletivo somado aos esforços
ponderados por fatores de rolagem, arfagem e guinada próprios da sua posição e
do seu sentido de rotação. Quando a soma excede a faixa disponível, a
alocação dá prioridade à rolagem e à arfagem sobre a guinada, preservando
apenas uma folga mínima para a guinada (\codigo{MOT\_YAW\_HEADROOM}), e
sinaliza a saturação ao controlador. Por fim, o empuxo desejado $t \in [0,1]$
de cada motor é convertido em comando $c$ pela inversa da curva de empuxo
$t = (1-e)\,c + e\,c^2$:
\begin{equation}
c = \frac{-(1-e) + \sqrt{(1-e)^2 + 4\,e\,t}}{2\,e},
\label{eq:linearizacao-empuxo}
\end{equation}
com $e$ dado por \codigo{MOT\_THST\_EXPO}, de modo que o esforço do controlador
seja aproximadamente linear em empuxo; o comando é então mapeado na faixa útil
do motor e convertido em largura de pulso.

\subsection{O \textit{framework} \texorpdfstring{\codigo{AC\_CustomControl}}{AC\_CustomControl}}
\label{sec:customcontrol}

O \codigo{AC\_CustomControl} é a biblioteca do ArduPilot destinada a
hospedar controladores de atitude alternativos sem alterar o restante do
firmware~\cite{ardupilot2026}. Sua função de atualização é chamada no laço
rápido logo depois do controlador de velocidade angular nativo e antes da
saída para os motores (Figura~\ref{fig:cascata}); ela recebe a mesma atitude
alvo, a mesma alimentação direta e as mesmas estimativas usadas pelo
controlador nativo e devolve os esforços de rolagem, arfagem e guinada,
que sobrescrevem os do nativo antes da alocação. A latência adicional é,
portanto, nula.

O \textit{framework} separa uma interface comum (\textit{frontend}) das
implementações (\textit{backends}). O parâmetro \codigo{CC\_TYPE} escolhe o
\textit{backend} --- 1 para um modelo vazio e 2 para um PID de exemplo, com
parâmetros próprios de prefixo \codigo{CC2\_} --- e \codigo{CC\_AXIS\_MASK}
escolhe, por máscara de bits, em quais eixos a saída do controlador
customizado substitui a do nativo. A ativação em voo é feita por uma chave do
rádio configurada com a função 109. Na volta ao controlador nativo, o
\textit{framework} zera os integradores do PID nativo e, com
\codigo{ATC\_RATE\_FF\_ENAB} habilitado, leva a atitude alvo e o alvo de
velocidade angular ao estado atual da aeronave e reinicializa os filtros dos
PIDs, para suavizar a transição; enquanto o \textit{backend} está ativo, os
integradores do nativo são zerados a cada laço nos eixos que ele controla. O
\textit{backend} de exemplo reproduz a arquitetura nativa de ângulo e
velocidade angular sem as suas proteções: o controlador de ângulo é só
proporcional, sem o controlador de raiz quadrada e sem a atenuação da guinada
acima de \ang{30}, e o PID de velocidade angular não recebe o sinal de
saturação da alocação. Os ganhos padrão são 90\,\% dos padrões compilados do
nativo; por isso serve como verificação do procedimento de ensaio, e não como
controlador alternativo. Os controladores propostos neste trabalho são
implementados como dois novos
\textit{backends} (Capítulo~\ref{cap:projeto}).

% ----------------------------------------------------------------------
\section{Técnicas de controle}
\label{sec:tecnicas}

\subsection{Controle PID em tempo discreto}

O controlador PID combina uma ação proporcional ao erro, uma ação integral,
que elimina o erro em regime diante de perturbações constantes, e uma ação
derivativa, que antecipa a tendência do erro~\cite{astrom2008,ogata2010}. Na
implementação digital, três cuidados são essenciais: a derivada deve ser
filtrada, porque a diferenciação amplifica o ruído de medição; o integrador
deve ser protegido contra \textit{windup} quando o atuador satura; e a
discretização deve preservar o comportamento contínuo na faixa de frequências
de interesse, o que exige taxa de amostragem bem acima da banda passante da
malha~\cite{franklin1998,wang2020}. A Equação~\eqref{eq:pid-taxa} mostra como o
ArduPilot trata esses três pontos, e acrescenta um termo de alimentação
direta que reduz a carga sobre a realimentação no rastreamento.

\subsection{Realimentação de estados e regulador linear quadrático}

Na realimentação de estados, a lei de controle $\vetor{u} = -\matriz{K}\vetor{x}$
aloca os autovalores de malha fechada $\matriz{A} - \matriz{B}\matriz{K}$
sempre que o par $(\matriz{A}, \matriz{B})$ é
controlável~\cite{ogata2010,fadali2009}.
No regulador linear quadrático (LQR), em vez de especificar os autovalores, o
projetista especifica o compromisso entre desvio dos estados e esforço de
controle por meio do custo
\begin{equation}
\mathcal{J} = \sum_{k=0}^{\infty}
\left(\vetor{x}_k^{\mathsf{T}}\matriz{Q}\,\vetor{x}_k
    + \vetor{u}_k^{\mathsf{T}}\matriz{R}_u\,\vetor{u}_k\right),
\qquad \matriz{Q} \succeq 0,\ \matriz{R}_u \succ 0,
\label{eq:custo-lqr}
\end{equation}
cujo mínimo, para o sistema discreto
$\vetor{x}_{k+1} = \matriz{A}_d\vetor{x}_k + \matriz{B}_d\vetor{u}_k$, é
atingido pelo ganho~\cite{franklin1998,datta2003}
\begin{equation}
\matriz{K} = \left(\matriz{R}_u + \matriz{B}_d^{\mathsf{T}}\matriz{P}\matriz{B}_d\right)^{-1}
             \matriz{B}_d^{\mathsf{T}}\matriz{P}\matriz{A}_d,
\label{eq:ganho-lqr}
\end{equation}
em que $\matriz{P}$ é a solução estabilizante, simétrica e semidefinida
positiva, da equação algébrica de Riccati discreta
\begin{equation}
\matriz{P} = \matriz{A}_d^{\mathsf{T}}\matriz{P}\matriz{A}_d
 - \matriz{A}_d^{\mathsf{T}}\matriz{P}\matriz{B}_d
   \left(\matriz{R}_u + \matriz{B}_d^{\mathsf{T}}\matriz{P}\matriz{B}_d\right)^{-1}
   \matriz{B}_d^{\mathsf{T}}\matriz{P}\matriz{A}_d + \matriz{Q}.
\label{eq:riccati}
\end{equation}
Essa solução existe e é única quando o par $(\matriz{A}_d, \matriz{B}_d)$ é
estabilizável e o par $(\matriz{A}_d, \matriz{Q}^{1/2})$ é
detectável~\cite{datta2003}; nos eixos deste trabalho, as duas condições
valem, porque a planta é controlável e $\matriz{Q}$ é diagonal positiva. Uma
escolha inicial usual para as ponderações é normalizar cada estado e cada
entrada pelo seu valor máximo admissível (regra de Bryson; ver
\citeonline{hespanha2009}), refinando-as a seguir pela resposta
de malha fechada. Com realimentação de todos os estados, o LQR de tempo
contínuo tem margens de estabilidade garantidas em cada canal de entrada:
tolera aumento ilimitado ou redução de até 50\,\% do ganho e tem margem de
fase de ao menos \ang{60}~\cite{safonov1977}. Em tempo discreto, mesmo com
realimentação de todos os estados, as margens garantidas são menores e
dependem de $\matriz{R}_u$ e de $\matriz{B}_d^{\mathsf{T}}\matriz{P}\matriz{B}_d$~\cite{shaked1986};
com estimador, deixam de existir~\cite{doyle1978}; e os filtros e atrasos da
planta real as reduzem ainda mais. Por isso as margens são verificadas em
cada projeto.

A realimentação de estados pura não anula o erro em regime diante de uma
perturbação constante, que é o caso de $\vetor{\tau}_d$ em degrau. A solução
clássica~\cite{young1972} é aumentar o vetor de estados com a integral do erro de
rastreamento, $x_{I,k+1} = x_{I,k} + T_s\,(r_k - y_k)$, e projetar o LQR sobre
o sistema aumentado, chamado LQI (\textit{linear quadratic integral}), o que
preserva a formulação de projeto ao custo de um
estado a mais por eixo e da necessidade de proteção contra \textit{windup}.

\subsection{Observador de estado estendido}

Uma alternativa à ação integral é estimar a perturbação e cancelá-la. No
controle por rejeição ativa de perturbações (ADRC), sistematizado em inglês
por \citeonline{han2009}, todo o efeito que separa a planta de um modelo nominal
simples --- perturbações externas, erros de modelo e acoplamentos --- é
agrupado em uma \emph{perturbação total} $f$, tratada como estado adicional e
estimada por um observador de estado estendido (ESO). Para um eixo de atitude
com a velocidade angular medida, a dinâmica se escreve
$\dot p = b_0\,u + f$, com $b_0$ a estimativa do ganho de controle, e um ESO
linear de segunda ordem é~\cite{gao2003}
\begin{equation}
\dot{\hat p} = \hat f + b_0\,u + l_1\,(p - \hat p), \qquad
\dot{\hat f} = l_2\,(p - \hat p),
\label{eq:eso}
\end{equation}
com $l_1$ e $l_2$ escolhidos para alocar os polos do observador. A lei de
controle $u = (u_0 - \hat f)/b_0$ reduz o eixo, dentro da banda do observador,
ao duplo integrador nominal, sobre o qual $u_0$ é projetado. O compromisso
central está na banda do observador: quanto maior, mais rápida a estimativa da
perturbação e maior a sensibilidade ao ruído do giroscópio.

\subsection{Filtro de Kalman e controle LQG}
\label{sec:lqg}

A lei $\vetor{u} = -\matriz{K}\vetor{x}$ supõe todos os estados medidos sem
erro. Quando parte deles não é medida, ou é medida com ruído e por um caminho
com dinâmica própria, o estado é estimado a partir de um modelo. Seja o
sistema discreto
\begin{equation}
\vetor{x}_{k+1} = \matriz{A}_d\vetor{x}_k + \matriz{B}_d\vetor{u}_k
                + \matriz{G}\,\vetor{w}_k,
\qquad
\vetor{y}_k = \matriz{C}\vetor{x}_k + \vetor{v}_k,
\label{eq:modelo-estocastico}
\end{equation}
com $\vetor{w}_k$ e $\vetor{v}_k$ ruídos brancos independentes, de
covariâncias $\matriz{W}$ e $\matriz{V}$. O filtro de
Kalman~\cite{kalman1960b} corrige a predição do modelo pela inovação, a
diferença entre a medida e a sua predição:
\begin{equation}
\hat{\vetor{x}}_{k|k} = \hat{\vetor{x}}_{k|k-1}
    + \matriz{L}\left(\vetor{y}_k - \matriz{C}\,\hat{\vetor{x}}_{k|k-1}\right),
\qquad
\hat{\vetor{x}}_{k+1|k} = \matriz{A}_d\hat{\vetor{x}}_{k|k} + \matriz{B}_d\vetor{u}_k.
\label{eq:kalman}
\end{equation}
Em regime, o ganho $\matriz{L}$ é constante e sai de uma equação de Riccati
com a forma de~\eqref{eq:riccati}, com $\matriz{A}_d^{\mathsf{T}}$ e
$\matriz{C}^{\mathsf{T}}$ no lugar de $\matriz{A}_d$ e $\matriz{B}_d$ e com as
covariâncias $\matriz{G}\matriz{W}\matriz{G}^{\mathsf{T}}$ e $\matriz{V}$ no
lugar das ponderações: é a dualidade entre o filtro e o
regulador~\cite{kalman1960b}. Como no ESO, o vetor de estados pode ser
estendido com a perturbação, modelada como passeio aleatório, que o filtro
passa a estimar; e pode ser estendido também com a dinâmica do caminho de
medida --- filtros, atrasos e polarizações dos sensores ---, que o ESO
de~\eqref{eq:eso} ignora. É o que faz o controlador LQG deste trabalho
(Capítulo~\ref{cap:projeto}).

O controlador LQG (\textit{linear quadratic Gaussian}) aplica o ganho do LQR
ao estado estimado, $\vetor{u}_k = -\matriz{K}\hat{\vetor{x}}_{k|k}$, e trata
o projeto como um problema estocástico de regulação~\cite{athans1971}. Pelo
princípio da separação~\cite{joseph1961}~\cite[seção~8.2]{anderson1989}, os
autovalores da malha fechada são a união dos autovalores da realimentação de
estados, $\matriz{A}_d - \matriz{B}_d\matriz{K}$, com os do erro de
estimação, e $\matriz{K}$ e $\matriz{L}$ podem ser projetados separadamente.
A separação vale para a estabilidade nominal, mas não para a robustez: as
margens garantidas do LQR com realimentação de todos os estados não se
transferem ao LQG, cujas margens dependem do projeto e precisam ser
verificadas~\cite{doyle1978}. Uma forma de recuperá-las é a recuperação da
malha (LTR, \textit{loop transfer recovery}): aumentando a intensidade do
ruído fictício de processo na entrada da planta, a função de transferência de
malha na entrada da planta tende à do LQR com realimentação de todos os
estados~\cite{doyle1979,doyle1981}~\cite[seção~8.4]{anderson1989}. A
recuperação exige planta de fase mínima e, em tempo discreto e com atraso na
medida, pode ser incompleta; o seu preço é um estimador mais rápido, que deixa
passar mais ruído de medida para o esforço de controle
(Capítulo~\ref{cap:projeto}).

\subsection{Limite de desempenho: a integral de Bode}
\label{sec:bode}

Seja $L(z)$ o ganho de malha aberta de um eixo, que inclui o controlador e a
planta $P(z)$, e $S(z) = 1/[1 + L(z)]$ a função de sensibilidade. Um torque
de perturbação que entra na planta chega ao ângulo por $P(z)\,S(z)$, a
sensibilidade à perturbação na entrada da planta~\cite{landau2011}, e não por
$S(z)$: é essa a grandeza que as varreduras de torque deste trabalho estimam
(Seção~\ref{sec:identificacao}). Para $L(z)$ estritamente própria e sem polos
fora do círculo unitário, como nos eixos deste trabalho, a sensibilidade em
tempo discreto obedece a~\cite{sung1988,mohtadi1990}
\begin{equation}
\frac{1}{\pi}\int_0^{\pi}\ln\left|S\!\left(\mathrm{e}^{\mathrm{j}\vartheta}\right)\right|
\mathrm{d}\vartheta = 0,
\qquad \vartheta = \omega\,T_s,
\label{eq:bode}
\end{equation}
versão discreta da integral de sensibilidade de
\citeonline{bode1945}; com polos instáveis em malha aberta, o lado direito
passa a ser positivo~\cite{freudenberg1985}. A atenuação da perturbação numa
faixa de frequências ($|S| < 1$) é paga com amplificação em outra
($|S| > 1$), em áreas iguais de $\ln|S|$: é o efeito conhecido como colchão
d'água, que nenhuma lei de controle linear e invariante no tempo contorna;
\citeonline{iglesias2001} estende a restrição a sistemas lineares variantes no
tempo. Em tempo discreto, a integral cobre uma faixa finita, de 0 à
frequência de Nyquist, e a amplificação não pode ser diluída numa banda
arbitrariamente larga, como no caso contínuo; com o laço a \SI{400}{\hertz}
e o filtro de \SI{20}{\hertz}, o compromisso fica mais apertado.
Como a planta $P$ é a mesma para todos os controladores, a diferença entre os
$\ln|P\,S|$ medidos com dois controladores é a diferença entre os seus
$\ln|S|$, e a comparação das varreduras de torque herda a
restrição~\eqref{eq:bode}.

% ----------------------------------------------------------------------
\section{Modelos de perturbação}
\label{sec:perturbacoes}

\subsection{Perturbações de torque}

Um vento uniforme aplicado ao centro de massa produz força, mas não momento:
afeta a posição e, apenas indiretamente, a atitude. Em uma aeronave real, o
vento também gera momentos --- pelo afastamento entre o centro de pressão e o
centro de massa, pelo \textit{blade flapping} e pela incidência desigual nos
rotores --- e são esses momentos que exercitam a malha de atitude. Para
avaliar a malha de atitude de forma direta e controlada, empregam-se
perturbações aditivas de torque, isto é, o termo $\vetor{\tau}_d$ de
\eqref{eq:rotacao-euler}, em quatro formas: o \emph{pulso retangular}, cujo
impulso angular é conhecido de antemão e que mede a recuperação após um
deslocamento; o \emph{degrau sustentado}, que expõe o erro em regime e a
suficiência da ação integral; a \emph{varredura em frequência}, que
caracteriza, ao longo de uma faixa de frequências, a resposta do ângulo ao
torque de perturbação, isto é, a sensibilidade à perturbação na entrada da
planta (Seção~\ref{sec:bode}); e o \emph{torque estocástico}, um processo de
Gauss-Markov de primeira ordem em cada eixo --- o processo de
Ornstein-Uhlenbeck~\cite{uhlenbeck1930} ---, que representa uma perturbação
contínua de banda limitada e é avaliado pelo erro quadrático médio. Com
desvio-padrão $\sigma$ e tempo de correlação $\tau$, o processo admite
discretização exata~\cite{gillespie1996}:
\begin{equation}
x_{k+1} = \mathrm{e}^{-\Delta t/\tau}\,x_k
        + \sigma\sqrt{1 - \mathrm{e}^{-2\Delta t/\tau}}\;w_k,
\qquad w_k \sim \mathcal{N}(0,1),
\label{eq:gauss-markov-discreto}
\end{equation}
que preserva a variância $\sigma^2$ para qualquer passo $\Delta t$.

Para que os ensaios sejam reprodutíveis, o ruído branco $w_k$ deve ser gerado
de forma determinística a partir de uma semente. Neste trabalho, ele é
obtido de um gerador congruencial linear multiplicativo,
$x_{k+1} = 16807\,x_k \bmod (2^{31}-1)$, o padrão mínimo de
\citeonline{park1988}, calculado sem estouro em aritmética de 32 bits pelo
método de \citeonline{schrage1979}, cujas amostras uniformes são convertidas em amostras
gaussianas pela transformação de \citeonline{box1958}. A mesma transformação,
na variante polar de \citeonline{marsaglia1964}, que dispensa seno e cosseno,
é a que o SITL usa (\codigo{Aircraft::rand\_normal}) na turbulência nativa do
vento e no ruído dos sensores do modelo de avião; no multirrotor, o ruído dos
sensores inerciais simulados é uniforme.

\subsection{Vento}

O vento que atua sobre a aeronave é usualmente decomposto em uma parcela
constante e uma parcela de rajadas, modelada como processo
estocástico~\cite{beard2012}. No SITL, o vento uniforme age quase só sobre a
posição (Seção~\ref{sec:sitl}), e nenhum ensaio de atitude deste trabalho
depende dele. Os modelos de vento do injetor --- o vento médio, a rajada
discreta da norma MIL-F-8785C~\cite{mil8785c} e a turbulência de Dryden, com
a sua discretização --- estão no Apêndice~\ref{ap:vento}.

% ----------------------------------------------------------------------
\section{Identificação de sistemas no domínio da frequência}
\label{sec:identificacao}

A identificação no domínio da frequência estima a resposta em frequência de
um sistema a partir de registros de entrada e saída obtidos sob uma excitação
de banda larga~\cite{ljung1999,pintelon2012,tischler2012}. A excitação mais usada em
aeronaves é a varredura senoidal de frequência crescente (\textit{chirp}). O
modo \textit{System Identification} do ArduCopter (modo 25) gera uma
varredura exponencial, implementada em
\codigo{AP\_Math/chirp.cpp}~\cite{ardupilot2026}, cuja frequência instantânea
é
\begin{equation}
\omega(t) = \omega_{\min}\,\exp\!\left(B\,\frac{t - t_c}{T_r - t_c}\right),
\qquad B = \ln\frac{\omega_{\max}}{\omega_{\min}},
\label{eq:chirp}
\end{equation}
precedida por um intervalo de duração $t_c$ em frequência constante
$\omega_{\min}$ e modulada por janelas de cosseno no início e no fim do
registro de duração $T_r$, de modo que a excitação comece e termine suave.
A varredura exponencial dedica o mesmo tempo a cada década, o que favorece as
baixas frequências.

A resposta em frequência é estimada pela razão entre a densidade espectral
cruzada da entrada $u$ com a saída $y$ e a densidade espectral da
entrada~\cite{carter1973,tischler2012},
\begin{equation}
\hat H(f) = \frac{\hat G_{uy}(f)}{\hat G_{uu}(f)},
\qquad
\hat\gamma^2_{uy}(f) = \frac{|\hat G_{uy}(f)|^2}{\hat G_{uu}(f)\,\hat G_{yy}(f)},
\label{eq:frf}
\end{equation}
com as densidades estimadas pelo método de \citeonline{welch1967}, que divide
o registro em segmentos sobrepostos, aplica uma janela a cada um e faz a
média dos periodogramas, trocando resolução em frequência por menor
variância. A coerência $\hat\gamma^2_{uy} \in [0, 1]$ mede a fração da saída
linearmente explicada pela entrada: valores baixos indicam ruído,
não linearidade ou excitação insuficiente naquela frequência, e
\citeonline{tischler2012} recomendam usar apenas os pontos de coerência
elevada (da ordem de 0,6 ou mais) no ajuste de modelos paramétricos. O ajuste
de uma função de transferência a $\hat H(f)$ por mínimos quadrados fornece o
modelo identificado, comparado ao modelo analítico
\eqref{eq:duplo-integrador}~\cite{ljung1999,pintelon2012}.

% ----------------------------------------------------------------------
\section{Simulação \textit{software-in-the-loop}}
\label{sec:sitl}

No SITL do ArduPilot, o próprio código do firmware é compilado para o
computador de desenvolvimento, com uma camada de abstração de
\textit{hardware} (HAL) que substitui os sensores e os atuadores por um modelo
físico do veículo~\cite{ardupilotdev}. O código de estimação e de controle é
o mesmo que roda na controladora de voo; o que muda é a origem das medidas e
o destino dos comandos de PWM. O relógio é simulado e avança junto com o
modelo físico, de modo que os resultados não dependem da carga do computador,
e a comunicação com o veículo usa o protocolo MAVLink, como em um voo real.

\subsection{Modelo de multirrotor}

O modelo físico de multirrotor do SITL, implementado em
\codigo{SIM\_Frame.cpp} e \codigo{SIM\_Motor.cpp}, na pasta
\codigo{libraries/SITL} do firmware, calcula, a cada passo, o empuxo e o
torque de cada motor a partir do seu comando $c \in [0, 1]$:
\begin{equation}
T_i = T_{\max}\left[(1-e)\,c_i + e\,c_i^2\right], \qquad
Q_i = 0{,}05\,D\,T_i\,c_i,
\label{eq:motor-sitl}
\end{equation}
em que $e$ é o expoente da curva de empuxo e $D$ é o parâmetro de tamanho do
quadro; a variação dos comandos é limitada em taxa. O torque de reação cresce,
portanto, com o comando, e a razão $\kappa = Q_i/T_i$ de
\eqref{eq:alocacao} vale $0{,}05\,D\,c$. Quando o tensor de inércia não é
informado, o simulador o estima supondo metade da massa distribuída em um anel
de raio $D/2$:
$J_{xx} = J_{yy} = m\,(D/2)^2/4$ e $J_{zz} = m\,(D/2)^2/2$.

A dinâmica de rotação é integrada no referencial do corpo como
\begin{equation}
\dot{\vetor{\omega}} = \matriz{J}^{-1}\left(\vetor{\tau} + \vetor{\tau}_{md}\right)
                     - a_r\,\vetor{\omega} + \vetor{a}_{\mathrm{tw}},
\label{eq:rotacao-sitl}
\end{equation}
em que $\vetor{\tau}_{md}$ é o torque do arrasto de momento dos rotores ---
cada rotor, ao se deslocar com a rotação do corpo, sofre uma força contrária à
sua velocidade ---, $a_r$ é um amortecimento rotacional linear explícito, de
constante de tempo da ordem de \SI{17}{\second}, e
$\vetor{a}_{\mathrm{tw}}$ é uma aceleração angular aditiva, ajustável pelos
parâmetros \codigo{SIM\_TWIST\_*}, que realiza exatamente o termo
$\matriz{J}^{-1}\vetor{\tau}_d$ de \eqref{eq:rotacao-euler}. É o arrasto de
momento que domina o amortecimento: com ele, a constante de tempo da
velocidade angular cai para cerca de \SI{0.5}{\second}
(Seção~\ref{sec:movimento}), ainda lenta perto da banda do controle. O termo
giroscópico $\vetor{\omega}\times\matriz{J}\vetor{\omega}$ \emph{não} é
modelado. Para uma aeronave com $J_{xx} = J_{yy}$, ele só afetaria rolagem e
arfagem, por meio do produto de duas velocidades angulares, e é desprezível
perto do pairado; a consequência é que a planta simulada está mais próxima do
modelo linear \eqref{eq:duplo-integrador} do que uma aeronave real estaria. O
vento (\codigo{SIM\_WIND\_*}) entra no modelo pela velocidade relativa do ar,
por dois caminhos: o arrasto de corpo, aplicado no centro de massa e sem
momento, e o arrasto de momento de cada rotor, aplicado na posição do rotor e
proporcional à raiz do seu empuxo. Com empuxos iguais, os momentos deste
último se cancelam, e com empuxos diferentes sobra um momento pequeno: o vento
move a aeronave, mas praticamente não perturba a atitude. Essa é a razão pela
qual os ensaios de rejeição de perturbação na malha de atitude usam
perturbações de torque.

\subsection{Sensores inerciais simulados}
\label{sec:sensores-sitl}

Com \codigo{INS\_FAST\_SAMPLE} ligado, como nos ensaios, o giroscópio
simulado é amostrado oito vezes por leitura, e a leitura entregue ao firmware
é a média das oito sub-amostras (\codigo{AP\_InertialSensor\_SITL.cpp}). A
cada sub-amostra, o SITL soma à velocidade angular verdadeira um ruído
uniforme de $\pm\SI{0.04}{\degree\per\second}$, sempre presente, e, com os
motores ligados, um segundo ruído uniforme, de amplitude igual a
\codigo{SIM\_GYR1\_RND} vezes o acelerador. Com a média, um ruído uniforme de
amplitude $A$ por sub-amostra chega ao firmware com desvio-padrão
$A/\sqrt{3 \cdot 8}$: o ruído mínimo fica da ordem de
\SI{0.01}{\degree\per\second}, e o giroscópio dos ensaios com
\codigo{SIM\_GYR1\_RND} nulo é, assim, praticamente sem ruído. Há ainda uma
deriva lenta, de \SI{0.05}{\degree\per\second} por minuto
(\codigo{SIM\_DRIFT\_SPEED}), que o EKF3 estima como polarização.

O filtro passa-baixas de \SI{20}{\hertz} (\codigo{INS\_GYRO\_FILTER}) só se
aplica à velocidade angular entregue aos controladores. Os incrementos de
ângulo que alimentam o estimador de atitude são integrados do giroscópio
antes do filtro (\codigo{AP\_InertialSensor\_Backend.cpp}): em horizontes
curtos, o ângulo estimado acompanha a integral do giroscópio sem o atraso do
filtro, mas com o ruído integrado, e chega ao controlador com um laço de
atraso (Seção~\ref{sec:instrumentacao}). O ruído, a média das sub-amostras e
o ângulo calculado sem o filtro definem o caminho de medida que o estimador do
controlador LQG modela (Capítulo~\ref{cap:projeto}).

\subsection{Aeronave simulada}

Os ensaios usam o modelo de quadrirrotor X padrão do SITL, cujos parâmetros
estão na Tabela~\ref{tab:aeronave}. Com eles, a alavanca de rolagem e arfagem
é $\ell_{xy} = \SI{0.247}{\metre}$, enquanto a razão entre torque de reação e
empuxo no pairado é $\kappa = \SI{0.0096}{\metre}$, nas condições de
referência do modelo. Como o torque de reação
cresce com o comando, a razão que vale para \emph{variações} de empuxo é
maior, cerca de \SI{0.0165}{\metre} (Seção~\ref{sec:motor-sitl}); ainda
assim, um mesmo padrão de diferença de empuxo produz um torque de guinada
cerca de 15 vezes menor que o de rolagem e, como $J_{zz} = 2J_{xx}$, uma
aceleração angular cerca de 30 vezes menor. A assimetria reaparece nas
acelerações máximas que o AutoTune registrou para o modelador de referência:
\SI{115}{\degree\per\second\squared} em guinada, contra cerca de
\SI{5800}{\degree\per\second\squared} em rolagem e arfagem
(Tabela~\ref{tab:ganhos}). A razão, de 50, não coincide com a de autoridade,
porque o AutoTune grava a aceleração atingida nos seus ensaios de ângulo, que
depende também dos ganhos, e não a máxima disponível.

\begin{table}[htbp]
\centering
\caption{Parâmetros do quadrirrotor X simulado.}
\label{tab:aeronave}
\begin{tabular}{lll}
\toprule
Grandeza & Valor & Origem \\
\midrule
Massa $m$                                   & \SI{3.0}{\kilogram} & \codigo{SIM\_FRM\_MASS} \\
Braço $\ell$ (centro ao motor)              & \SI{0.35}{\metre}   & \codigo{SIM\_FRM\_DIAG\_SZ}\textsuperscript{(a)} \\
Alavanca de rolagem e arfagem $\ell_{xy}$   & \SI{0.247}{\metre}  & $\ell/\sqrt{2}$ \\
Inércia $J_{xx} = J_{yy}$                   & \SI{0.02297}{\kilogram\metre\squared} & estimada pelo SITL \\
Inércia $J_{zz}$                            & \SI{0.04594}{\kilogram\metre\squared} & estimada pelo SITL \\
Empuxo máximo por motor $T_{\max}$          & \SI{18.86}{\newton} & calculado pelo SITL \\
Empuxo por motor no pairado                 & \SI{7.35}{\newton}, com $c = 0{,}551$ & $mg/4$ \\
Razão empuxo/peso                           & 2,56 & --- \\
Expoente da curva de empuxo $e$             & 0,65 & \codigo{SIM\_FRM\_EXPO} \\
Razão $\kappa$ no pairado                   & \SI{0.0096}{\metre} & $0{,}05\,D\,c$ \\
Taxa do modelo físico                       & \SI{1200}{\hertz} & \codigo{SIM\_RATE\_HZ} \\
Taxa do laço de controle                    & \SI{400}{\hertz} & \codigo{SCHED\_LOOP\_RATE} \\
\bottomrule
\end{tabular}


\fonte{Parâmetros gravados nos ensaios e código-fonte do SITL
\cite{ardupilot2026}.}
\nota{\raggedright (a) Apesar do nome (``diagonal''), o SITL posiciona cada motor a essa
distância do \emph{centro}; a distância entre motores opostos é, por isso,
\SI{0.70}{\metre}.}
\end{table}
